Este trabalho apresenta um estudo de caso sobre a identificação e a validação do modelo dinâmico não linear para o modo multirrotor no protótipo DH (Drone Híbrido), um veículo aéreo não tripulado (VANT) VTOL. É nesse modo que ocorrem a decolagem e o pouso verticais, fases essenciais para a operação autônoma, que demandam um modelo dinâmico fiel para o projeto de controle. Para obter esse modelo, emprega-se uma formulação de caixa cinza, na qual a estrutura parte das equações de corpo rígido, enquanto os parâmetros físicos da aeronave, como momentos de inércia, coeficientes de empuxo e de contratorque dos rotores e arrastos angulares, são estimados a partir de dados de voos reais. Como o multirrotor é instável e só voa estabilizado por um controlador, identifica-se a planta em malha aberta a partir de respostas em malha fechada, contornando a divergência numérica e a correlação entre entrada e estados. Esse processo segue a M4V, metodologia de identificação no domínio do tempo. As manobras de excitação, do tipo \textit{doublet}, são aplicadas a um eixo de cada vez, para isolar a dinâmica de cada canal. Os dados de voo resultantes passam por sincronização, remoção de viés e filtragem. Sobre esses dados, os parâmetros são estimados em duas fases: o método do erro de equação fornece a estimativa inicial algébrica, refinada pelo método do erro de saída, que minimiza a diferença entre a resposta simulada e as medições. O modelo identificado é validado em janelas de voo distintas das de estimação, com movimentos compostos que excitam os três eixos simultaneamente, apresentando boa concordância entre as respostas simuladas e medidas. O eixo de guinada mostra-se o mais difícil de identificar, pela baixa autoridade de controle e pela elevada inércia nesse eixo. Na sequência, esse modelo é linearizado em torno do voo pairado e, sobre a representação linear resultante, projeta-se o controle em cascata com controladores PID e suas variações, sintonizados por alocação de polos. A malha fechada é simulada nas plantas linear e não linear, mapeando a faixa de manobras em que a linearização permanece representativa. Como contribuição, o trabalho entrega a cadeia completa do dado bruto de voo ao controlador simulado e um modelo validado que serve de base para pilotos automáticos de decolagem vertical, mais aderente aos dados de voo que o obtido apenas com medidas geométricas e ensaios de bancada.

\keywords{Identificação de Sistemas. VTOL. Método do Erro de Saída. Controle PID em cascata.}
